\documentclass[12pt]{article}

\usepackage[a4paper, margin=2.5cm]{geometry}
\usepackage{amsmath}
\usepackage{amssymb}
\usepackage{amsthm}
\usepackage[colorlinks=true, linkcolor=blue, citecolor=blue, urlcolor=blue]{hyperref}
\usepackage{booktabs}
\usepackage{natbib}

\newtheorem{theorem}{Theorem}
\newtheorem{proposition}[theorem]{Proposition}
\newtheorem{corollary}[theorem]{Corollary}
\newtheorem{lemma}[theorem]{Lemma}
\theoremstyle{definition}
\newtheorem{definition}[theorem]{Definition}
\newtheorem{hypothesis}[theorem]{Hypothesis}
\theoremstyle{remark}
\newtheorem{remark}[theorem]{Remark}

\newcommand{\Gcal}{\mathcal{G}}
\newcommand{\farcs}{\mbox{$\,''\!$}}

\title{Schwarzschild-Target Ansatz on a Heuristic Closure-Field Graph Model\\[6pt]
\large Spherically Symmetric Construction and Weak-Field Targets (Conditional on Paper XL and Paper XXXVI; staticity deferred to a future continuum-side definition; substrate bridge to the VFD event poset not supplied)}
\author{%
  Lee Smart\\[4pt]
  \textit{Institute of Vibrational Field Dynamics}\\[2pt]
  \texttt{contact@vibrationalfielddynamics.org}\\[2pt]
  \texttt{@vfd\_org}%
}
\date{April 2026}

\begin{document}

\maketitle

\noindent\textbf{Status:} Preprint (not peer-reviewed). Conditional on Paper XL continuum-limit targets.

\begin{abstract}
This paper sets up a spherically symmetric closure-field ansatz on a discrete model of the kind introduced in~\citep{SmartXXVII} (staticity is deferred to a future continuum-side definition: on the discrete vertex set there is no time coordinate, only Paper~XXIII's gauge-fixed birth-angle order, and the substrate bridge needed to translate ``static'' into that setting is identified as absent in Paper~XL), and records the conditions under which one would seek a derivation reducing the ansatz to the Schwarzschild--de~Sitter weak-field metric: Paper~XL's Conjecture~T2 (or its Proposition~``T2 under BI-sharp and halving'' with those hypotheses verified), a precise formulation and proof of XL's T3 programme target, realisation of XL's T4 programme conjecture (including the absent substrate-bridging projection and the tensorial upgrade), a separate pointed local-Euclidean regime near the source (a local asymptotic hypothesis, not the global claim ``$\mathbb{R}^3$ asymptotically flat at infinity''), a Kottler / Birkhoff reduction of the vacuum Einstein-with-$\Lambda$ sector (so that the Cascade~$\Lambda$ value of~\citep{SmartXXXVI} becomes the $-\Lambda r^2/3$ term), import of Paper~XXXVI's F8 and Cascade~$\Lambda$ theorems with their additional model assumptions, and an additional scalar normalisation fixing the Poisson coupling $\lambda = 4\pi G$. The paper provides a heuristic ansatz and inventory, not a derivation; no component of the chain is proved here.

Heuristic targets (all conditional on the above; labelled explicitly as heuristic, not proved):
\begin{itemize}
\item (\emph{Structural, heuristic}) A spherically symmetric ansatz on the closure-field model, with staticity deferred to a future continuum-side definition and a source-strength parameter $M_{\mathrm{src}}$ entering the scalar target equation; separate later hypotheses are needed to identify $M_{\mathrm{src}}$ with a normalised physical mass $M$ and, in the exterior-metric target below, with the Kottler mass parameter $m_{\mathrm{Kott}}$.
\item (\emph{Heuristic target}) Newtonian limit $\phi_{\mathrm{Newton}}(r) \to -GM/r$ at leading order in $GM/r$, corresponding to $g_{00} \to -(1 - 2GM/r)$ under the additional assumption that the full continuum tensorial limit agrees with Einstein gravity at first post-Newtonian order.
\item (\emph{Heuristic target}) Recovery of the first-PN GR values $\gamma_{\mathrm{PPN}} = \beta_{\mathrm{PPN}} = 1$ \emph{if} the full first-PN metric emerges from the (as-yet-unconstructed) tensorial continuum theory.
\item (\emph{Structural, heuristic}) Candidate black-hole-horizon analogue in the weak-field Schwarzschild benchmark ($\Lambda \to 0$ comparison limit), where the continuum radius $r = 2GM$ is the Schwarzschild horizon; under T3-SS's Kottler target metric the actual target horizon is a root of $1 - 2GM/r - \Lambda r^2/3 = 0$ rather than $r = 2GM$ exactly. The proposed discrete pre-image is a ``closure-coherence boundary'' (no such notion is currently established in the VFD observer framework), labelled as a construction target rather than a proved correspondence.
\end{itemize}
The paper does not derive the first-PN metric or the PPN coefficients. It does not construct the continuum limit (that is the subject of~\citep{SmartXL}). It is the programme's first concrete Schwarzschild-target paper in the GR branch, intended to name the target form and to record most of the required upstream inputs from Papers~XL, XXVII, and XXXVI (XXIX is relevant only to the speculative horizon-analogue section), together with additional local-Euclidean, shell-incidence, source-limit, and source-to-exterior matching assumptions made explicit below.
\end{abstract}

%==================================================================
\section{Introduction}\label{sec:intro}
%==================================================================

Papers XXIII--XXVIII of the VFD programme develop event-order, observer, separation-function, curvature-indicator, and graph-Poisson constructions on the event-poset framework derived from the dual 600-cell~\citep{SmartXXIII, SmartXXIV, SmartXXV, SmartXXVI, SmartXXVII, SmartXXVIII}; several of those results are on toy models or under explicit conditions, and the continuum limit is deferred. Paper~XL~\citep{SmartXL} organises the continuum-limit questions into targets T1--T4 and proves conditional versions of T1 weak and T1 strong under explicit hypotheses, gives T2 as a conditional corollary (under BI-sharp + halving) plus a separate unconditional conjecture, and states T3 as a programme target (Lorentzian objects and convergence notion not yet defined) and T4 as a programme conjecture / integration target (requiring among other things a substrate-bridging $H_4 \to D_4$ projection not currently available anywhere in the programme). The present paper is the programme's first concrete Schwarzschild-target paper: it sets up a spherically symmetric heuristic ansatz on a closure-field model (with staticity deferred to a future continuum-side definition, since the discrete vertex set has no time coordinate) and inventories what the continuum limit would have to produce for it to reduce to Schwarzschild--de~Sitter at leading weak-field order, only after an additional decompactification / local-Euclidean step (not provided by~\citep{SmartXL} itself) and a Kottler / Birkhoff reduction of the vacuum Einstein-with-$\Lambda$ sector (likewise not supplied by the programme). No solution, conditional or otherwise, is derived here.

\subsection{Scope}
\begin{itemize}
\item Set up the symmetry-reduced ansatz for a spherically symmetric closure-field configuration (staticity deferred to a future continuum-side definition).
\item Write the radial form of the graph-Poisson-style equation under that ansatz, with ``$\approx$'' understood as pointing to an unestablished continuum limit rather than as a proved discrete identity.
\item Inventory, target by target, what the XL/XXXVI programme would have to deliver for the ansatz to reduce to Schwarzschild--de~Sitter at leading weak-field order.
\item List the first-PN observables (redshift, light deflection, perihelion precession, Shapiro delay) that would follow under such a reduction, without deriving their coefficients from the present discrete construction.
\end{itemize}
\textbf{What this paper does not do}: it does not prove the continuum limit (that is the subject of Paper~XL); it does not derive the PN coefficients $\beta$ or $\gamma$ from the discrete ansatz; it does not supply the substrate-bridging projection that Paper~XL identifies as missing; it does not construct rotating, charged, or time-dependent solutions; and it does not establish any ``coherence boundary'' notion beyond naming it as a construction target.

%==================================================================
\section{The Symmetry-Reduced Ansatz}\label{sec:ansatz}
%==================================================================

\subsection{Spherical ansatz (staticity deferred)}

We work on a finite vertex set $X$ equipped with a graph-distance $d$; the reader may think of $X_n$ from Paper~XL's Schl\"afli refinement sequence with $d = \delta_n$ (cf.~\citep{SmartXL}), but the present ansatz does not require, and does not supply, the substrate-bridging projection identified as absent in Paper~XL (\citealt{SmartXL}, T4 prerequisite). Choose a distinguished origin vertex $o \in X$. A closure field $\Phi : X \to \mathbb{R}$ is \emph{spherically symmetric} (with staticity deferred; see SSS-2 below) if:
\begin{itemize}
\item[(SSS-1)] $\Phi(x)$ depends only on the graph-distance $d(x, o)$, not on angular direction.
\item[(SSS-2)] Schematically, $\Phi$ is intended to be ``time-independent'' in the yet-to-be-constructed Lorentzian continuum limit. On the discrete side $X$, however, no time coordinate exists: Paper~XXIII~\citep{SmartXXIII} supplies only a gauge-fixed strict birth-angle order on events, not a metric time coordinate, and the substrate bridge from $X$ to that event-order structure is identified as absent in Paper~XL~\citep{SmartXL}. The present paper therefore does not \emph{define} SSS-2 as a condition on $X$; we retain the label only as a reminder that a future substrate bridge would eventually need to specify what ``static'' means relative to the event-order parameter, and the remaining analysis uses only SSS-1.
\end{itemize}
Accordingly, the operative discrete symmetry below is SSS-1 alone; SSS-2 is retained only as a heading for the future definitional task, and no equation in this paper is derived from it.

Denote the radial profile by $\phi(r) = \Phi(x)$ where $r = d(x, o)$. Throughout this paper, the symbol $r$ is used in two distinct senses: (i) as a discrete shell index $r \in \mathbb{Z}_{\geq 0}$ equal to the graph-distance from the origin (in equations~\eqref{eq:radial-recurrence}--\eqref{eq:origin} and the surrounding discrete analysis), and (ii) as a continuum radial coordinate $r \in \mathbb{R}_{>0}$ (in the T1-SS target, T3-SS, and the PN observables). The identification between the two senses requires an additional step --- the identification of the discrete shell index with a continuum radial coordinate with uniform spacing --- which the present paper assumes informally and does not establish.

\subsection{The discrete Poisson-analogue equation}

Paper~XXVII~\citep{SmartXXVII} defines a scalar graph-Poisson model $\Delta_\Gcal u = S - \bar S$ on the Hasse graph of the event poset --- note the subtraction of the source mean $\bar S$, which is the finite-connected-graph compatibility condition forced by $\sum_x (\Delta_\Gcal \Phi)(x) = 0$ --- and explicitly disclaims any Einstein equation, tensorial stress-energy, continuum limit, and coupling constant. \emph{As a continuum-limit target, and not as a solvable finite-graph equation}, we consider the schematic source equation
\begin{equation}\label{eq:discrete-poisson}
(\Delta \Phi)(x) \;\sim\; \lambda \cdot T(x),
\end{equation}
where $\Delta$ is a graph Laplacian on $X$ (the specific graph and its bridge to the Schl\"afli refinement sequence of~\citep{SmartXL} are not supplied by the present programme; cf.~\citep{SmartXL} T4 prerequisite), $\lambda$ is a positive scalar that the programme eventually hopes to identify with the Newtonian Poisson coupling $4\pi G$ (the value that makes the continuum scalar equation $\nabla^2 \phi = 4\pi G \rho$ yield $\phi = -GM/r$ for a point source of strength $M$). Fixing that identification would require an additional normalisation argument not supplied in the programme; in particular, Paper~XXXVI's F8 performs action-coefficient matching using imported couplings and structural model choices (\citealt{SmartXXXVI}), and does not by itself determine $\lambda$ in terms of Newton's constant. The Einstein-field-equation normalisation $8\pi G$, which is sometimes invoked in the programme in connection with the absent tensorial theory, is a different coupling from the scalar Poisson $\lambda$. $T(x)$ is an idealised point source at the origin for the vacuum-point-mass ansatz:
\begin{equation}
T(x) = M_{\mathrm{src}} \cdot \delta_o(x), \qquad \delta_o(x) = \begin{cases} 1 & x = o \\ 0 & x \ne o. \end{cases}
\end{equation}
The symbol $M_{\mathrm{src}}$ denotes the source strength of the scalar target equation. It is \emph{a priori} distinct from (i) a normalised physical mass $M$ (to be identified with $M_{\mathrm{src}}$ only by the separate normalisation argument under T1-SS below) and (ii) the Kottler mass parameter $m_{\mathrm{Kott}}$ (to be identified with $M$ only by the separate source-to-exterior matching assumption under T3-SS). In displays below we suppress these subscripts once the identifications $M_{\mathrm{src}} = M = m_{\mathrm{Kott}}$ have been made explicit as hypotheses; the reader should interpret each use of $M$ in the light of the nearest of those hypotheses.
On a finite connected graph, $\Delta \Phi = \lambda M \delta_o$ as written has no solution (the right-hand side has nonzero sum while the left has zero sum); a legitimate finite-graph formulation would have to use XXVII's mean-subtracted form $\Delta \Phi = \lambda(T - \bar T)$, a punctured / infinite-volume setting on which $\delta_o$ becomes the Laplacian of the Green's function, or an explicit boundary-flux specification. The present paper does not adopt one of these; it uses~\eqref{eq:discrete-poisson} only as a continuum-limit target equation, with ``$\sim$'' intended to point toward $\nabla^2 \phi = \lambda M \delta^3(\mathbf{r})$ in the continuum regime where the finite-graph obstruction is absent. The symbol $\lambda$ is used here to avoid collision with the BI-bi-Lipschitz constant $\kappa$ of~\citep{SmartXL}; no identification of $\lambda$ with $8\pi G$ is derived here.

We now introduce an auxiliary continuum-inspired radial template. This is \emph{not} a reduction of the graph Laplacian on Paper~XL's Schl\"afli refinement sequence: such a reduction would require a defined shell geometry, a uniform-spacing identification of the discrete shell index with a continuum radial coordinate, and a spherical-symmetry structure on the graph itself, none of which is supplied by~\citep{SmartXL}. Under SSS-1 alone, and anticipating the continuum limit rather than deriving it, we \emph{import} the 3D Euclidean radial-Laplacian template as a heuristic target equation. In finite-difference form with unit spacing this reads
\begin{equation}\label{eq:radial-recurrence}
\phi(r+1) - 2\phi(r) + \phi(r-1) + \frac{1}{r}\bigl(\phi(r+1) - \phi(r-1)\bigr) \;\approx\; 0 \qquad (r \geq 1),
\end{equation}
the standard central-difference discretisation of $\phi''(r) + (2/r) \phi'(r) = 0$, together with a separate origin/source condition which (as noted above) cannot literally be the pointwise identity $(\Delta \Phi)(o) = \lambda M$ on a finite connected graph but should be read as a schematic origin input
\begin{equation}\label{eq:origin}
(\Delta \Phi)(o) \;\sim\; \lambda M \qquad \text{(schematic; not a finite-graph equation),}
\end{equation}
corresponding in the continuum to $\nabla^2 \phi = 0$ away from the source and $\nabla^2 \phi = \lambda M \delta^3(\mathbf{r})$ at the source; under a pointed local-Euclidean limit, local Green-function normalisation data on a regime $r \ll (6GM/\Lambda)^{1/3}$ (so that the $\Lambda$-corrections of T3-SS are locally negligible), and flat shell-volume scaling (each a separate assumption not established by the discrete construction above), the choice $\lambda = 4\pi G$ would make the continuum Green function $\phi = -GM/r$ the target bulk solution. The ``$\approx$'' in~\eqref{eq:radial-recurrence} is not a proved discrete identity: no quantitative error term is given, and translating the graph Laplacian on a generic Schl\"afli refinement into a radial recurrence requires additional hypotheses (a spherical symmetry of the graph itself, uniform spacing along radial shells, identification of the graph-distance shell index with a continuum radial coordinate) that are not established for~\citep{SmartXL}'s refinement sequence. We retain equations~\eqref{eq:radial-recurrence}--\eqref{eq:origin} as a heuristic target system, not as a theorem.

%==================================================================
\section{Continuum Limit and the Schwarzschild Solution}\label{sec:continuum}
%==================================================================

\subsection{Newtonian-limit target}

\textbf{T1-SS (heuristic target, not a proved proposition).} If
\begin{itemize}
\item Paper~XL's Conjecture~T2 holds unconditionally, or, more modestly, it is enough that its Proposition~``T2 under BI-sharp and halving'' holds with those hypotheses verified (GH-convergence of a connected tail of the Schl\"afli refinement sequence to $(S^3, d_{\mathrm{round}})$; \citealt{SmartXL}),
\item Paper~XL's programme target T3 is made precise and established (a Lorentzian discrete structure $\{(X_n^{\mathrm{Lor}}, \delta_n^{\mathrm{Lor}})\}$ and a convergence notion --- both undefined in~\citep{SmartXL} --- supplying a Minkowski-type limit object),
\item Paper~XL's programme conjecture T4 is realised (tensorial upgrade of~\citep{SmartXXVII}'s scalar graph-Poisson model, substrate-bridging projection from $H_4$ to $D_4$, and the continuum limit of the tensorial equation),
\item Paper~XXXVI's F8 action-coefficient matching and the associated Einstein--Hilbert normalisation are imported with their model-hypothesis load (\citealt{SmartXXXVI}, F8), and
\item a separate scalar normalisation argument fixes the Poisson coupling to the Newtonian value $\lambda = 4\pi G$ and identifies the source parameter $M$ with physical mass,
\item a source-limit prescription is given --- we take, as a specific choice among several possibilities mentioned around~\eqref{eq:discrete-poisson}, the mean-subtracted finite-graph equation $\Delta \Phi = M_{\mathrm{src}}(\delta_o - 1/|X|)$ at each level, equipped with the gauge $\Phi(o) = 0$ (or any other pointwise normalisation chosen once and kept fixed across all levels) to select a unique finite-graph representative, together with a scaling of $M_{\mathrm{src}}$ in $n$ and a source-limit scheme (vanishing background subtraction plus concentration at the origin) such that the rescaled $\Phi$ converges to a continuum delta-source Poisson problem on the pointed Euclidean space of the next bullet; the \emph{particular} Green-kernel normalisation of that continuum problem is fixed by the gauge/scaling choice, and is \emph{not} assumed up-front to equal $-GM/r$ ---
\item the radial-reduction assumptions of Section~\ref{sec:ansatz} are met: spherical symmetry of the graph itself, a uniform-spacing identification of the discrete shell index with a continuum radial coordinate, flat 3-space shell-volume scaling, and a shell-to-shell incidence / edge-count asymptotic that matches the central-difference radial Laplacian $\phi'' + (2/r)\phi'$ --- none of which is established for Paper~XL's Schl\"afli refinement sequence,
\item an explicit blow-up / rescaling limit around the source is given, of standard pointed Gromov--Hausdorff form: for $n \geq N^*$ with the connected-tail hypothesis of Paper~XL, rescale the metric $\delta_n$ by the level-$n$ covering radius $\varepsilon_n \to 0$ to obtain a pointed metric space $(X_n, o, \varepsilon_n^{-1}\delta_n)$, and assume that this pointed sequence converges in pointed Gromov--Hausdorff distance to $(\mathbb{R}^3, 0, d_{\mathrm{Euclid}})$. This is a \emph{local} asymptotic-flatness condition on a small neighbourhood of the source, not a global claim, and it is not established for Paper~XL's Schl\"afli refinement sequence in any round of the present programme,
\end{itemize}
then one may ask whether an additional pointed local-Euclidean limit near the source --- boundary data chosen so that the local Euclidean delta-source Poisson problem on $\mathbb{R}^3$ is well-posed --- yields the Newtonian gravitational potential
\begin{equation}
\phi_{\mathrm{Newton}}(r) \;=\; -\frac{GM}{r}
\end{equation}
as the continuum limit of the radial profile $\phi(r)$ taken by \emph{importing} the 3D Euclidean radial-Laplacian template of~\eqref{eq:radial-recurrence}--\eqref{eq:origin} rather than by reducing the graph Laplacian of~\citep{SmartXXVII} through~\citep{SmartXL}'s refinement sequence (a reduction that is not carried out in the present paper). Here $G$ is Newton's constant, $M$ is the physical mass identified with the source-strength parameter $M_{\mathrm{src}}$ \emph{only under the separate normalisation hypothesis above} (outside that hypothesis $M_{\mathrm{src}} \neq M$ in general). Each bullet above is an open item: of the first four, Paper~XL's T2 is an unconditional conjecture (with a conditional proposition available under BI-sharp and halving), XL's T3 is a programme target not yet precisely formulated, XL's T4 is a programme conjecture that is not yet a precise mathematical conjecture (and requires the absent $H_4 \to D_4$ substrate-bridging projection), and Paper~XXXVI's F8 is proved as a coefficient-matching theorem --- but F8 does not by itself fix the Newtonian scalar normalisation $\lambda = 4\pi G$ that the present target requires. The fifth bullet is not supplied by any paper in the programme. The pointed local-Euclidean step near the source is a local asymptotic hypothesis compatible with a nonzero $\Lambda$ elsewhere: it does \emph{not} require the global spacetime to be asymptotically flat, and in particular it is consistent with the Kottler asymptotics required by T3-SS below, provided $r \ll (6GM/\Lambda)^{1/3}$ as discussed there.

\begin{remark}[Why T1-SS is not a proposition]
The present paper does not derive the continuum limit (that is the XL programme, explicitly unfinished), does not supply the substrate bridge (XL T4 prerequisite, absent), and does not establish a normalisation argument fixing $\lambda = 4\pi G$ (in particular, F8's action-coefficient matching does not by itself determine the scalar Poisson coupling). The result also implicitly assumes flat-3-space shell-volume scaling and uniform radial-shell spacing in the graph Laplacian --- neither is established for the Schl\"afli refinement sequence. Each of these is either an open conjecture, an open programme target, or a missing argument. Treating T1-SS as a theorem would mean upgrading at least one of them; we therefore label it a heuristic target, not a proposition.
\end{remark}

\subsection{First post-Newtonian target}

\textbf{T2-SS (heuristic target, not a proved proposition).} If the target T1-SS above is realised, and \emph{additionally} the continuum tensorial theory produced by XL's T4 agrees with Einstein gravity at first post-Newtonian order (this is a separate open assumption: the PN coefficients are not determined by action-normalisation matching alone), then the effective metric functions admit the GR first-PN form
\begin{equation}
g_{00}(r) \;=\; -1 + \frac{2GM}{r} + \mathcal{O}\!\left(\frac{(GM)^2}{r^2}\right), \qquad \beta_{\mathrm{PPN}} = \gamma_{\mathrm{PPN}} = 1.
\end{equation}

\begin{remark}[Why $\beta=\gamma=1$ is not derived]
Target~T1-SS would recover only the Newtonian potential $\phi_{\mathrm{Newton}} = -GM/r$, hence only the $2GM/r$ piece of $g_{00}$ via the weak-field identification $g_{00} = -(1 + 2\phi)$. That does not fix (i) the $(GM/r)^2$ coefficient in $g_{00}$ (which controls $\beta$) and (ii) the $GM/r$ coefficient in $g_{ij}$ (which controls $\gamma$); both are data of the full tensorial theory, not of a scalar graph-Poisson limit. Paper~XXXVI's F8 is action-coefficient matching with imported couplings and structural model choices; it does not determine the PN coefficients of the Schwarzschild weak-field solution. T2-SS is therefore a target form, not a theorem.
\end{remark}

\subsection{Schwarzschild--de~Sitter extension target}

\textbf{T3-SS (heuristic target, not a proved proposition).} If, in addition to T1-SS and T2-SS,
\begin{itemize}
\item Paper~XXXVI's conditional Cascade~$\Lambda$-Theorem is imported with \emph{all} of its additional model assumptions (the shell-amplitude rule, the product-closure rule, and the residue-to-$\Lambda$ identification; \citealt{SmartXXXVI}, and \citealt{SmartXL} preserves exactly this dependency load),
\item a separate continuum-side notion of \emph{static spherical symmetry} is supplied (this paper defers staticity; SSS-2 is undefined on the discrete side per Section~\ref{sec:ansatz}, so the continuum ``static spherically symmetric'' condition needed below is an extra prerequisite, not an output of the present construction), and
\item the absent continuum vacuum sector additionally satisfies the Einstein-with-$\Lambda$ equations $R_{\mu\nu} - \tfrac{1}{2} R g_{\mu\nu} + \Lambda g_{\mu\nu} = 0$ strongly enough that, under the supplied static-spherical-symmetry condition, each connected static spherically symmetric vacuum exterior region is locally isometric to a Kottler (Schwarzschild--de~Sitter) member, under the usual 4D Lorentzian vacuum-Einstein-with-$\Lambda$ regularity hypotheses required for the Kottler / Birkhoff classification,
\end{itemize}
then the relevant imported GR exterior benchmark is the Kottler family, parametrised by a mass parameter $m_{\mathrm{Kott}}$ (with dimensions of mass; $G m_{\mathrm{Kott}}$ has dimensions of length, consistent with the metric coefficient $2 G m_{\mathrm{Kott}} / r$ below):
\begin{equation}
\mathrm{d}s^2 \;=\; -\left(1 - \frac{2 G m_{\mathrm{Kott}}}{r} - \frac{\Lambda r^2}{3}\right) \mathrm{d}t^2 + \left(1 - \frac{2 G m_{\mathrm{Kott}}}{r} - \frac{\Lambda r^2}{3}\right)^{-1} \mathrm{d}r^2 + r^2 \mathrm{d}\Omega^2,
\end{equation}
and a further \emph{source-to-exterior matching assumption} is required to identify $m_{\mathrm{Kott}} = M$ with the physical mass of T1-SS (itself identified with the scalar source strength $M_{\mathrm{src}}$ only under T1-SS's separate normalisation hypothesis); Birkhoff's theorem alone fixes the Kottler family but not this identification. Subject to that matching,
\begin{equation}
g_{00}(r) \;=\; -\left(1 - \frac{2GM}{r} - \frac{\Lambda r^2}{3}\right),
\end{equation}
with $\Lambda \ell_P^2 = 2\varphi^{-583}$ coming from the conditional Cascade~$\Lambda$-Theorem. Importing the Cascade~$\Lambda$-Theorem alone gives only the numerical value of $\Lambda$; the $-\Lambda r^2/3$ term itself requires the Kottler reduction just named, which is an additional assumption not established anywhere in the programme. The $\Lambda r^2/3$ term is the asymptotic de~Sitter contribution; for fixed source mass $M$, the $\Lambda$-term is negligible relative to the $2GM/r$ Schwarzschild term only in the regime $\Lambda r^3 \ll 6GM$, equivalently $r \ll (6GM/\Lambda)^{1/3}$ (a mass-dependent crossover), not in the mass-independent regime $r \ll \Lambda^{-1/2}$. The shell-amplitude, product-closure, and residue-to-$\Lambda$ assumptions of the cited theorem are not derived here.

%==================================================================
\section{Post-Newtonian Observables}\label{sec:pn-obs}
%==================================================================

If targets T1-SS, T2-SS, and T3-SS were realised, the target exterior metric would be Kottler; and \emph{provided every radius sampled by a given observable} (source radius, observer radius, and impact parameter / perihelion / closest-approach radius, as appropriate) lies in the local regime $\Lambda r^3 \ll 6GM$, the $\Lambda$ corrections are negligible relative to the Schwarzschild $2GM/r$ terms, so the standard Schwarzschild first-PN formulas are the intended local benchmarks for that observable. We list these leading Schwarzschild first-PN observables in geometrized units ($c = 1$, consistent with the $r_s = 2GM$ and $\omega/\omega_\infty$ formulas used elsewhere in this paper), with $\Lambda$ corrections dropped in the local regime, to make the target of the programme explicit. None of the following formulas is derived from the discrete ansatz in the present paper; they are the GR values that the programme's eventual Schwarzschild-target derivation would have to reproduce.

\begin{description}
\item[\textbf{PN1 --- Gravitational redshift}] For a photon emitted at radius $r$ and received by a static observer at larger radius $r_0 > r$ inside the local-Euclidean regime, $\omega(r)/\omega(r_0) = \sqrt{(1-2GM/r_0)/(1-2GM/r)} \approx 1 + GM(1/r - 1/r_0)$ (greater than $1$: the received frequency is smaller than the emitted frequency, so light from deeper in the potential well is redshifted). The asymptotic-infinity Schwarzschild expression $\omega(r)/\omega_\infty = (1 - 2GM/r)^{-1/2}$ is the standard asymptotically-flat benchmark to which this local formula reduces as $r_0 \to \infty$; that benchmark is given here only for comparison, not as the statement that the programme adopts globally. \emph{GR experimental status}: Pound--Rebka~\citep{PoundRebka1960} and GPS time-dilation tests consistent with GR at $\lesssim 10^{-10}$.
\item[\textbf{PN2 --- Light deflection}] Total bending angle $\Delta\phi = 4GM_\odot/r_\odot$ (in geometrized units; $= 1\farcs75$ in physical units). \emph{GR experimental status}: Eddington~\citep{Eddington1919}; VLBI~\citep{Fomalont2009} at $\sim 10^{-4}$.
\item[\textbf{PN3 --- Perihelion precession}] Mercury: $\Delta\varpi = 6\pi GM_\odot / [a(1-e^2)]$ per orbit (in geometrized units; $= 43''\,\mathrm{per\,century}$ in physical units); consistent with observation.
\item[\textbf{PN4 --- Shapiro time delay}] Round-trip radar delay past the Sun: $\Delta t = 4GM_\odot \ln(4 r_1 r_2 / b^2)$ (in geometrized units). \emph{GR experimental status}: Cassini~\citep{Bertotti2003} at $\gamma_{\mathrm{PPN}} - 1 = (2.1 \pm 2.3) \times 10^{-5}$.
\end{description}

\begin{remark}[No VFD-specific prediction is derived here]
Under the target form, the first-PN observables agree with GR by construction. Any VFD-specific deviation would first appear at second PN order or via higher-shell $\varphi^{-n}$ corrections to the closure coupling; the present paper computes neither, and no VFD-specific numerical prediction is produced.
\end{remark}

%==================================================================
\section{A Future Horizon-Analogue Target (Sharply Demoted)}\label{sec:horizon}

\begin{remark}[Scope of this section]
The present paper is advertised as a weak-field target-setting paper; its discrete side does not yet supply staticity, a time coordinate, or a causal structure (cf.\ SSS-2 at Section~\ref{sec:ansatz} and the deferred-staticity notice in the abstract and introduction). ``Event horizon'' is, by contrast, a global causal notion of smooth Lorentzian geometry, and any VFD correspondence must wait on a construction that is explicitly absent from the current programme. We therefore use the section below only to \emph{name} a candidate future target and to record what would be needed to state it precisely; we do not propose a discrete horizon theorem.
\end{remark}

%==================================================================

In the $\Lambda \to 0$ Schwarzschild benchmark, the continuum metric has an event horizon at $r_s = 2GM$; under T3-SS's full Kottler target metric with $\Lambda > 0$ (and $9 G^2 M^2 \Lambda < 1$, the usual sub-extremal condition), $1 - 2GM/r - \Lambda r^2/3 = 0$ has two positive roots: the smaller root $r_{\mathrm{bh}}$ is the black-hole horizon (reducing to $r_s = 2GM$ as $\Lambda \to 0$), and the larger root $r_{\mathrm{c}}$ is the cosmological horizon (tending to $\sqrt{3/\Lambda}$ as $M \to 0$). A natural VFD question is whether the discrete closure-field model admits a structural pre-image of the black-hole horizon (Schwarzschild $r_s = 2GM$ in the $\Lambda \to 0$ benchmark, or Kottler $r_{\mathrm{bh}}$ under T3-SS). T4-SS concerns this black-hole horizon; the cosmological horizon at $r_{\mathrm{c}}$ is a separate, globally distinct surface and is not addressed here. We \emph{name} such a candidate pre-image the ``closure-coherence boundary'' and record it as a construction target; we do not establish it.

\textbf{T4-SS (construction target, not a proved proposition).} Under the targets T1-SS--T3-SS, the Kottler black-hole horizon radius $r_{\mathrm{bh}}$ (the smaller positive root of $1 - 2GM/r - \Lambda r^2/3 = 0$, which reduces to $r_s = 2GM$ in the $\Lambda \to 0$ Schwarzschild benchmark) would be the continuum image of some closure-coherence structure attached to the discrete source configuration --- informally, a radius below which the VFD closure construction fails a future, as-yet-undefined closure-coherence criterion. Supplying this structure requires defining, on the discrete side, what a ``closure-coherence boundary'' actually is; no such notion is currently supplied by Paper~XXIX~\citep{SmartXXIX} or elsewhere in the programme. The nearest existing machinery in XXIX is its observer-admissible configurations and localised-probe observer construction, which is distinct from a horizon-type boundary; cf.~\citep{SmartXXIX}.

\begin{remark}[Relation to Paper XXIX]
Paper~XXIX supplies observer-admissible configurations and closure-basin constructions for localised probes; it does not supply a horizon analogue or a theorem about observer trajectories becoming non-extendable inward of a critical radius. The ``coherence boundary'' language of the present paper is therefore a naming convention for a future construction, not a citation to an existing XXIX result.
\end{remark}

%==================================================================
\section{Scope and Open Questions}\label{sec:scope}
%==================================================================

\textbf{What this paper sets up / records} (heuristic / target-setting only; nothing is proved unconditionally):
\begin{itemize}
\item A spherically symmetric closure-field ansatz on a finite graph with origin; on the discrete side the operative condition is SSS-1 alone, and SSS-2 is retained only as a label for a future substrate-bridging / continuum-static notion not defined on $X$.
\item A radial target equation~\eqref{eq:radial-recurrence} imported as a Euclidean radial-Laplacian template, not reduced from the graph Laplacian of~\citep{SmartXXVII}, with ``$\approx$'' acknowledged as unsupported.
\item A named list of target statements T1-SS--T4-SS, each explicitly contingent on a stack of unproved upstream items.
\item An experimental-status summary for the four first-PN observables that would have to emerge from a successful realisation of T1-SS--T3-SS.
\end{itemize}

\textbf{Open items}:
\begin{enumerate}
\item Paper~XL's T2 (unconditional version), T3 (not yet a precise conjecture), and T4 (programme conjecture / integration target), including the absent $H_4 \to D_4$ substrate-bridging projection.
\item A tensorial upgrade of~\citep{SmartXXVII}'s scalar model.
\item Paper~XXXVI's F8 action matching, upgraded from a matching scheme with imported couplings to a derivation of the discrete-to-continuum Newton's constant.
\item Paper~XXXVI's conditional Cascade~$\Lambda$-Theorem, imported with its additional assumptions still explicit (shell-amplitude rule, product-closure rule, residue-to-$\Lambda$ identification).
\item A source-limit prescription (introduced locally at T1-SS) replacing the unsolvable finite-graph point-source equation by a mean-subtracted, punctured / infinite-volume, or boundary-flux limit that converges to a continuum delta source.
\item A shell-to-shell incidence / edge-count asymptotic on the discrete substrate that matches the central-difference 3D radial-Laplacian template~\eqref{eq:radial-recurrence}.
\item An explicit blow-up / rescaling limit around the source converting Paper~XL's $S^3$ target into a pointed Euclidean $\mathbb{R}^3$ Green-function problem.
\item A source-to-exterior matching theorem identifying the Kottler mass parameter $m_{\mathrm{Kott}}$ with the (normalised, identified) physical mass $M$ (Birkhoff's theorem alone fixes the Kottler family but not this identification).
\item Definition of a ``closure-coherence boundary'' on the discrete side and a theorem linking it to the Schwarzschild horizon $r = 2GM$ (in the $\Lambda \to 0$ benchmark) or the Kottler black-hole horizon $r_{\mathrm{bh}}$, the smaller positive root of $1 - 2GM/r - \Lambda r^2/3 = 0$ (under T3-SS).
\item Derivation (not just target form) of the first-PN GR coefficients $\beta_{\mathrm{PPN}} = \gamma_{\mathrm{PPN}} = 1$ from the (absent) tensorial continuum theory.
\item Second-PN corrections and any VFD-specific deviation from GR.
\item Rotating (Kerr-analogue) solutions, interior (matter) solutions, and stellar structure.
\end{enumerate}

\subsection{Conclusion}

The present paper does not derive a Schwarzschild analogue. It sets up a spherically symmetric ansatz on a closure-field graph model and records staticity as a deferred continuum-side task, writes down a target radial equation, and specifies what a future derivation of the Schwarzschild--de~Sitter weak-field metric would require: open items from Papers~XL, XXVII, and XXXVI (XXIX is relevant only to the speculative horizon-analogue section), together with the additional assumptions introduced here (source-limit prescription, shell-incidence asymptotic, local blow-up / rescaling, and source-to-exterior matching). No numerical VFD-specific prediction is produced, and no first-PN coefficient is derived. The paper's role in the programme is to name the target form of the first concrete gravitational solution and to inventory the work remaining before any part of the target chain could be upgraded to a theorem.

%==================================================================
\bibliographystyle{plain}
\begin{thebibliography}{99}

\bibitem{SmartXXIII} L. Smart, \emph{Event Structure and the Emergence of Time}, Paper XXIII, VFD Programme, 2026.
\bibitem{SmartXXIV} L. Smart, \emph{Observer Kinematics from the Event Poset}, Paper XXIV, VFD Programme, 2026.
\bibitem{SmartXXV} L. Smart, \emph{Metric Emergence from Event-Order Geometry}, Paper XXV, VFD Programme, 2026.
\bibitem{SmartXXVI} L. Smart, \emph{Curvature from Non-Uniform Event-Order Geometry}, Paper XXVI, VFD Programme, 2026.
\bibitem{SmartXXVII} L. Smart, \emph{Dynamics and Field Equations from Event-Order Geometry}, Paper XXVII, VFD Programme, 2026.
\bibitem{SmartXXVIII} L. Smart, \emph{Quantum and Gravitational Structure from a Common Event-Order Geometry}, Paper XXVIII, VFD Programme, 2026.
\bibitem{SmartXXIX} L. Smart, \emph{Observer as Constraint Substructure}, Paper XXIX, VFD Programme, 2026.
\bibitem{SmartXXXVI} L. Smart, \emph{Cascade Foundations: The Eight Foundations F1--F8}, Paper XXXVI, VFD Programme, 2026.
\bibitem{SmartXL} L. Smart, \emph{Continuum Targets and Conditional Reductions for $D_4$ and the GR Limit}, Paper XL, VFD Programme, 2026.

\bibitem{PoundRebka1960} R.\,V. Pound and G.\,A. Rebka, Jr., \emph{Apparent Weight of Photons}, Physical Review Letters \textbf{4}, 337 (1960).

\bibitem{Eddington1919} F.\,W. Dyson, A.\,S. Eddington, and C.\,R. Davidson, \emph{A Determination of the Deflection of Light by the Sun's Gravitational Field}, Philosophical Transactions of the Royal Society A \textbf{220}, 291 (1920).

\bibitem{Fomalont2009} E. Fomalont, S. Kopeikin, G. Lanyi, and J. Benson, \emph{Progress in Measurements of the Gravitational Bending of Radio Waves Using the VLBA}, The Astrophysical Journal \textbf{699}, 1395 (2009).

\bibitem{Bertotti2003} B. Bertotti, L. Iess, and P. Tortora, \emph{A test of general relativity using radio links with the Cassini spacecraft}, Nature \textbf{425}, 374 (2003).

\end{thebibliography}

\end{document}
